% ============================================================================
% 第 5 章外壳 —— 只负责装配四问，正文内容全在子文件里。
%
% 每问两片是这套结构的核心：
%   *-分析与准备.tex：数据体检、口径落地、预处理、记号约定、可行性预判
%   *-建模与求解.tex：变量/目标/约束/方程、算法与实现、结果与证书、图表解读
% 这两片加起来是全文体量最大的部分（基线中单问的「建模与求解」可达全文最大一节）。
% ============================================================================
\section{模型的建立与求解}

\subsection{问题一模型的建立与求解}
% ============================================================================
% 问题一 · 分析与准备
%
% 【四问通用合同，只在问题一处完整写一次，其余同构】
% 本片要交代清楚「在动模型之前，把地基打成什么样」：
%   1. 数据体检：规模、量纲、缺失/异常、与题面口径是否一致（给数字）
%   2. 口径落地：本问采用哪条口径，边界/退化情形怎么处理
%   3. 预处理与记号：坐标系、单位、索引约定
%   4. 可行性预判：先量单次计算耗时，倒推可承受的样本量/网格规模
%      —— 这一步写进正文，是 repro 与 algorithm 两维的加分点
% 篇幅参考：每问 1.5-3 页。
% ============================================================================
\subsubsection{分析与准备}

\input{05.1-问题一-建模与求解.tex}

\subsection{问题二模型的建立与求解}
\input{05.2-问题二-分析与准备.tex}
\input{05.2-问题二-建模与求解.tex}

\subsection{问题三模型的建立与求解}
\input{05.3-问题三-分析与准备.tex}
% ============================================================================
% 问题三 · 建模与求解
%
% 【四问通用合同】本片是全文体量最大的部分，必须写透这五段：
%   1. 模型建立：变量、目标、约束、方程。推导只留关键结论，全过程移附录。
%   2. 求解算法：复杂度、分辨率预算、样本量及其来源（怎么定的，不是拍的）。
%   3. 主结果：数字全部走宏，禁止手抄。
%   4. 证书：置信下界 / 收敛判据 / 双路互证，至少一样。只给点估计的结论要降措辞。
%      最优性主张分级：global / local / best-found 三选一，无据不得写「最优」。
%   5. 图表解读：每张图配一段「从图上看出什么」，不是只把图放上去。
%
% 【搜索类问题的额外红线】网格步长必须由可行前沿的陡度决定：
%   前沿附近退化为单位步长，或把由别问解出的阈值点强制塞进候选集。
%   两家候选都栽在「网格把最优点跨过去」上，这是高频失分点。
% 篇幅参考：每问 3-6 页。
% ============================================================================
\subsubsection{模型的建立}

\subsubsection{求解算法与实现}

\subsubsection{结果与统计证书}

\subsubsection{结果讨论}


\subsection{问题四模型的建立与求解}
\input{05.4-问题四-分析与准备.tex}
% ============================================================================
% 问题四 · 建模与求解
%
% 【四问通用合同】本片是全文体量最大的部分，必须写透：
%   1. 步骤式推导，每步公式链完整（密度契约见下）
%   2. 求解算法 + 参数表（含「选取依据」列）
%   3. 结果表 + 结果图（图必须进正文，未引用即 G5 FAIL）
%   4. 统计证书：置信下界 / 收敛判据 / 双路互证，至少一样
%   5. 结果讨论：数字 + 比较 + 机理，且显式引用本节的表与图
%   6. 结果叙事合同：现象—原因—意义；每个数字能回指 qN.json、aggregate.json 与 ledger。
%
% 【搜索类问题的额外红线】网格步长必须由可行前沿的陡度决定：
%   前沿附近退化为单位步长，或把由别问解出的阈值点强制塞进候选集。
%   两家候选都栽在「网格把最优点跨过去」上，这是高频失分点。
% 篇幅参考：每问 3-6 页；写长了再拆片。
% ============================================================================
\subsubsection{模型的建立}

% 开篇引文：一句话说清本节要建什么模型、从哪个角度切入。

% ---- 步骤式推导（这是全文体量的主要来源，也是评委判断建模深度的地方）----
% 【密度契约】按「步骤 1、2、3……」逐步推导，每个步骤：
%   - 先用一句话说清这一步要解决什么
%   - 再给出公式链（每步不少于 6 个编号公式；少于这个数通常意味着把推导跳过了）
%   - 每个新符号出现时当场解释，不要攒到符号表
%   - 步骤之间要有承接句，说明为什么下一步是这一步的必然延伸
% 推导的完整过程留在正文，只有冗长的中间代数才移附录——反过来做会把正文掏空。

\textbf{步骤 1：……}

\textbf{步骤 2：……}

\textbf{步骤 3：……}

\textbf{步骤 N：求解算法}

% 算法流程 + 参数表。注意参数表必须有「选取依据」一列——
% 只给取值不给依据，等于承认参数是拍的。
\begin{table}[htbp]
  \centering
  \caption{算法参数设置}
  \begin{tabularx}{\linewidth}{llXX}
    \toprule
    参数 & 取值 & 含义 & 选取依据 \\
    \midrule
     &  &  &  \\
    \bottomrule
  \end{tabularx}
\end{table}

% 承上启下：本节建立了什么，下一节据此求解什么。

\subsubsection{模型的求解}

% 求解环境、规模、耗时；样本量/网格规模怎么定的（先量单次耗时 -> 倒推可行规模
% -> 检查精度能否分辨相邻档），这一段是 algorithm 与 repro 两维的共同加分点。

% ---- 结果表 ----
\begin{table}[htbp]
  \centering
  \caption{……结果}
  \label{tab:qfour-result}
  \begin{tabularx}{\linewidth}{llXX}
    \toprule
     &  &  &  \\
    \bottomrule
  \end{tabularx}
\end{table}

% ---- 结果图（位置先摆好；本问的关键结果必须至少有一张图进正文）----
% 【硬规则】生成的图未被正文引用 = G5 FAIL。图做出来却没进论文，论证链就是断的。
% 历史教训：某轮交付物生成了 7 张图，正文只引用 1 张，figures 维直接归零。
\begin{figure}[htbp]
  \centering
  \IfFileExists{../图/fig_qfour_main.pdf}{%
    \includegraphics[width=0.85\linewidth]{fig_qfour_main.pdf}%
  }{%
    \IfFileExists{fig_qfour_main.pdf}{%
      \includegraphics[width=0.85\linewidth]{fig_qfour_main.pdf}%
    }{%
      \fbox{\parbox{0.82\linewidth}{\centering 图位占位：生成并接入本问证据图后再交稿}}%
    }%
  }
  \caption{……}
  \label{fig:qfour-main}
\end{figure}

\subsubsection{结果与统计证书}

% 主结果数字全部走宏。每个概率型结论旁必须给置信下界与样本量。
% 最优性主张分级：global / local / best-found，无据不得写「最优」。

\subsubsection{结果讨论}

% 【固定句式，确保表与图都被真正使用】
% 由表 \ref{tab:qfour-result} 可知……；从图 \ref{fig:qfour-main} 可以看出……；
% 综合以上结果，……（给机理解释：为什么是这个数，物理/结构上怎么理解）。
% 结果描述三要素缺一不可：数字 + 与基准的比较 + 机理。

\subsubsection*{本问结论}
% 【硬条件，2026-09-04 评委反馈：关键结论要更醒目地集中呈现】
% 只写三行，加框呈现，让评委扫读时一眼看到：
%   1. 最终答案（数字走宏，加粗）
%   2. 证书（样本量 + 置信下界 + 判定；用 \xxxCert 宏一并带出）
%   3. 一句话适用条件 / 口径限制
% 有 tcolorbox / mdframed 时用更好看的框；下面是纯 LaTeX 的保底写法。
\begin{center}
\fbox{\parbox{0.92\linewidth}{%
  \textbf{最终答案：}…… \\
  \textbf{证\qquad 书：}…… \\
  \textbf{适用条件：}……%
}}
\end{center}


\subsection{四问结果汇总}
% 【硬条件，2026-09-04 评委反馈】全文一张汇总表，一行一问。
% 评委据此一眼定位全部结论；没有这张表，结论就散在几十页里等人去找。
\begin{table}[htbp]
  \centering
  \caption{四问结果汇总}
  \label{tab:summary}
  \begin{tabularx}{\linewidth}{lXXlXl}
    \toprule
    问题 & 模型类别 & 求解算法 & 最终答案 & 证书 & 章节 \\
    \midrule
    一 &  &  &  &  &  \\
    二 &  &  &  &  &  \\
    三 &  &  &  &  &  \\
    四 &  &  &  &  &  \\
    \bottomrule
  \end{tabularx}
\end{table}
